\section{Audit of app229; corrected bracket bound via gauge-invariant
AF region; honest status of the mass-gap proof}
\label{app:audit_app229_corrected}

This appendix does five things:
\begin{enumerate}
\item \textbf{Audits app229} (\S\ref{ssec:app230_audit}): identifies
two errors in the proof of the Clay mass-gap conjecture
(Theorem~\ref{thm:clay_mass_gap_app229}).
\item \textbf{Retracts} the mass-gap claim of app229
(Theorem~\ref{thm:app229_retraction}).
\item \textbf{Identifies} the correctable error (the bracket bound)
and the non-correctable error (the scaling/normalization) in
\S\ref{ssec:app230_audit}.
\item \textbf{Provides a corrected bracket bound}
(\S\ref{ssec:app230_fix}) using a gauge-invariant AF region
$U_{\mathrm{AF},j}^{\mathrm{corr}}=\{\|F_{A}\|_{L^{2}}\le
\eta_{\mathrm{AF}}\,2^{2j}/(j+1)\}$.
\item \textbf{Establishes the honest final status}
(\S\ref{ssec:app230_status}): the mass gap is reduced to a single
gauge-invariant concentration estimate
(Hypothesis~\ref{hyp:gauge_invariant_concentration}).
\end{enumerate}

\subsection{Audit of app229}
\label{ssec:app230_audit}

\subsubsection{Error 1: The bracket bound has an unabsorbed
$(j+1)^{1/2}$ factor}

\begin{theorem}[Error in Lemma~\ref{lem:bracket_bound} of app229]
\label{thm:app229_error1}
The proof of Lemma~\ref{lem:bracket_bound} (app229, lines 399--408)
claims the bound
\begin{equation*}
  \bigl|\langle F_{A},[X\wedge X]\rangle_{L^{2}}\bigr|
  \;\le\;C_{\mathrm{br}}\,\eta_{\mathrm{AF}}\,
  \|D_{A}X\|_{L^{2}}^{2},
\end{equation*}
with $C_{\mathrm{br}}>0$ independent of $j$. The proof derives
the intermediate bound
\begin{equation}
\label{eq:app229_bracket_intermediate}
  \bigl|\langle F_{A},[X\wedge X]\rangle_{L^{2}}\bigr|
  \;\le\;C\,\eta_{\mathrm{AF}}\,(j+1)^{1/2}\,
  \|D_{A}X\|_{L^{2}}^{2},
\end{equation}
and then claims the factor $(j+1)^{1/2}$ is absorbed by the
``refined'' bound $\|F_{A}\|_{L^{2}}\le
\eta_{\mathrm{AF}}\,2^{2j}/(j+1)$.

\textbf{The refined bound does not hold.} On the AF region
$U_{\mathrm{AF},j}=\{\|A\|_{L^{2}}\le\eta_{\mathrm{AF}}\,
2^{j}/(j+1)^{1/2}\}$, the curvature satisfies
$F_{A}=dA+\tfrac{1}{2}[A,A]$ with
\begin{align*}
  \|dA\|_{L^{2}}
  &\;\le\;2^{j}\,\|A\|_{L^{2}}
  \;\le\;\eta_{\mathrm{AF}}\,\frac{2^{2j}}{(j+1)^{1/2}},\\
  \|[A,A]\|_{L^{2}}
  &\;\le\;\|A\|_{L^{4}}^{2}
  \;\le\;c_{\mathrm{Sob}}\,\|A\|_{L^{2}}^{2}
  \;\le\;c_{\mathrm{Sob}}\,\eta_{\mathrm{AF}}^{2}\,
  \frac{2^{2j}}{j+1}.
\end{align*}
The \emph{linear} term $\|dA\|_{L^{2}}\le
\eta_{\mathrm{AF}}\,2^{2j}/(j+1)^{1/2}$ dominates and gives
\begin{equation}
\label{eq:fa_bound_correct}
  \|F_{A}\|_{L^{2}}
  \;\le\;\eta_{\mathrm{AF}}\,\frac{2^{2j}}{(j+1)^{1/2}}
  \;+\;c_{\mathrm{Sob}}\,\eta_{\mathrm{AF}}^{2}\,
  \frac{2^{2j}}{j+1}
  \;\le\;2\,\eta_{\mathrm{AF}}\,\frac{2^{2j}}{(j+1)^{1/2}},
\end{equation}
for $j\ge 0$ (the quadratic term is smaller by a factor
$\eta_{\mathrm{AF}}/(j+1)^{1/2}\le 1$). The claimed bound
$\|F_{A}\|_{L^{2}}\le\eta_{\mathrm{AF}}\,2^{2j}/(j+1)$ is
\textbf{stronger by a factor $(j+1)^{-1/2}$} and does \emph{not}
follow from the given estimates.

Consequently, the bracket bound
\eqref{eq:app229_bracket_intermediate} with the factor
$(j+1)^{1/2}$ is the best available, and the Hessian lower bound
becomes
\begin{equation}
\label{eq:hessian_lower_app229_correct}
  \mathrm{Hess}\,S_{\mathrm{YM},j}(A)(X,X)
  \;\ge\;\frac{1}{g_{\mathrm{YM}}^{2}}\,
  \bigl(1-C_{\mathrm{br}}\,\eta_{\mathrm{AF}}\,(j+1)^{1/2}\bigr)\,
  \|D_{A}X\|_{L^{2}}^{2},
\end{equation}
which is \emph{negative} for $j>1/(C_{\mathrm{br}}\eta_{
\mathrm{AF}})^{2}-1$. The uniform coercivity
(Theorem~\ref{thm:be_curvature_gauge_invariant} of app229) fails.
\end{theorem}

\begin{proof}[Proof of Theorem~\ref{thm:app229_error1}]
The computation is given in the statement. The key point is that
$F_{A}=dA+\tfrac{1}{2}[A,A]$ has a \emph{linear} term in $A$,
and $\|dA\|_{L^{2}}\le 2^{j}\|A\|_{L^{2}}$ (the scale-$j$ Poincar\'e
upper bound), which on the AF region gives
$\|dA\|_{L^{2}}\le\eta_{\mathrm{AF}}\,2^{2j}/(j+1)^{1/2}$, not
$\eta_{\mathrm{AF}}\,2^{2j}/(j+1)$. The ``refinement'' claimed in
app229 conflates the linear term (scale $j\to j+1$, factor $2^{j}$)
with the quadratic term (scale $j\to 2j$, factor
$\|A\|_{L^{2}}^{2}\sim 2^{2j}/(j+1)$), but the linear term dominates.
\end{proof}

\subsubsection{Error 2: The scaling/normalization of
$h_{\mathrm{OS}}^{(n)}$ versus $h_{j}^{\nu}$}

\begin{theorem}[Scaling inconsistency in the multiscale aggregation]
\label{thm:app229_error2}
The proof of Theorem~\ref{thm:clay_mass_gap_app229} (app229,
\S\ref{ssec:app229_mass_gap}) computes the multiscale Cheeger
constant as
\begin{equation}
\label{eq:app229_multiscale_claim}
  h_{\mathrm{OS}}^{(n)}
  \;\ge\;\frac{c_{0}}{L_{\mathrm{phys}}}\,
  \biggl(\sum_{j=0}^{J}\frac{j+1}{2^{2j}}\biggr)^{-1/2}
  \;\ge\;\frac{c_{0}'}{L_{\mathrm{phys}}},
\end{equation}
which is \emph{independent of the lattice spacing $a_{n}$}. The
spectral gap is then bounded as
$1-\lambda_{1}^{(n)}\ge(c_{0}')^{2}/(2\mathfrak{D}_{*}\,
L_{\mathrm{phys}}^{2})$, also independent of $a_{n}$, giving
$m_{n}=(1-\lambda_{1}^{(n)})/a_{n}\to\infty$.

\textbf{This is inconsistent with the physical mass gap.} The
continuum mass gap $m_{\mathrm{phys}}=\lim_{n}m_{n}$ should be
\emph{finite and positive}, not infinite. The error is a
confusion between the per-scale Cheeger constant $h_{j}^{\nu}$
(in physical units, with explicit $L_{\mathrm{phys}}$ factor)
and the OS transfer-matrix Cheeger constant $h_{\mathrm{OS}}^{(n)}$
(in lattice units, with $a_{n}$-dependent normalization).

The correct relation (cf.\ app226, Theorem 1.5;
app218, Theorem~\ref{thm:rp_cheeger}) is
\begin{equation}
\label{eq:correct_scaling}
  1-\lambda_{1}^{(n)}
  \;\ge\;\frac{(h_{\mathrm{RP}}^{(n)})^{2}}
       {2\,(\beta_{n}N|P_{t}^{(n)}|+h_{\mathrm{RP}}^{(n)})},
  \;\;\;\;
  h_{\mathrm{RP}}^{(n)}\sim c\,a_{n}^{1/2},
\end{equation}
where $\beta_{n}\sim\log(1/a_{n})$ (asymptotic freedom) and
$|P_{t}^{(n)}|\sim L_{\mathrm{phys}}^{3}/a_{n}^{3}$ (number of
time-like plaquettes). The denominator
$\beta_{n}N|P_{t}^{(n)}|\sim L_{\mathrm{phys}}^{3}\log(1/a_{n})/
a_{n}^{3}$ diverges as $a_{n}\to 0$, and the correct scaling gives
$1-\lambda_{1}^{(n)}\sim a_{n}$, hence
$m_{n}\sim c>0$ (finite).
\end{theorem}

\begin{proof}[Proof of Theorem~\ref{thm:app229_error2}]
The app229 proof replaces the denominator
$\beta_{n}N|P_{t}^{(n)}|$ of the Cheeger--Lawler--Sokal bound
(app218, Theorem~\ref{thm:rp_cheeger}) with
$\mathfrak{D}_{*}^{\mathrm{eff}}$, claiming
$\mathfrak{D}_{*}^{\mathrm{eff}}\le\mathfrak{D}_{*}$ is
independent of $n$. However:
\begin{itemize}
\item The diffusivity $\mathfrak{D}_{*}$ of app226
(Theorem~\ref{thm:D_from_wms_corr}) is the WMS Sobolev constant,
which is indeed independent of $n$ in \emph{physical} units.
\item The Cheeger--Lawler--Sokal denominator
$\beta_{n}N|P_{t}^{(n)}|$ is in \emph{lattice} units and scales
as $a_{n}^{-3}\log(1/a_{n})$.
\item The conversion between physical and lattice units introduces
a factor $a_{n}^{-2}$ (from the OS inner product normalization),
which app229 omits.
\end{itemize}
The corrected bound is
$1-\lambda_{1}^{(n)}\ge(h_{\mathrm{OS}}^{(n)})^{2}/(2\mathfrak{D}_{
\mathrm{phys}}\,a_{n}^{-2}\cdot L_{\mathrm{phys}}^{2})$, where
$\mathfrak{D}_{\mathrm{phys}}$ is the physical-unit diffusivity.
With $h_{\mathrm{OS}}^{(n)}\sim c_{0}'/L_{\mathrm{phys}}$
(physical units), this gives
$1-\lambda_{1}^{(n)}\ge(c_{0}')^{2}a_{n}^{2}/(2\mathfrak{D}_{
\mathrm{phys}}L_{\mathrm{phys}}^{4})$, hence
$m_{n}\ge(c_{0}')^{2}a_{n}/(2\mathfrak{D}_{\mathrm{phys}}\,
L_{\mathrm{phys}}^{4})\to 0$.

\emph{The mass gap still vanishes} unless $h_{\mathrm{OS}}^{(n)}$
scales as $a_{n}^{-1/2}$ in physical units (i.e.\ $a_{n}^{1/2}$
in lattice units), which requires the per-scale Cheeger bound
$h_{j}^{\nu}\ge c\,2^{2j}/(j+1)$ (the $\epsilon=1/2$ case of
app226), not $h_{j}^{\nu}\ge c\,2^{j}/(j+1)^{1/2}$.
\end{proof}

\subsection{Retraction of the app229 mass-gap claim}
\label{ssec:app229_retraction}

\begin{theorem}[Retraction]
\label{thm:app229_retraction}
The Clay Millennium Prize mass-gap conjecture
(Theorem~\ref{thm:clay_mass_gap_app229} of app229) is
\textbf{retracted}. The proof contains two errors
(Theorems~\ref{thm:app229_error1} and~\ref{thm:app229_error2})
that invalidate the conclusion.
\end{theorem}

\begin{remark}[What survives from app229]
\label{rem:app229_survives}
The following results of app229 are \emph{not} affected by the
retraction and remain valid:
\begin{enumerate}
\item \textbf{The gauge-invariant quotient measure $\nu_{j}$}
(Definition~\ref{def:gauge_invariant_measure}) is well-defined
and is the natural OS measure on $\mathcal{Q}_{j}$.
\item \textbf{The horizontal Sobolev inequality}
(Lemma~\ref{lem:horizontal_sobolev}) is correctly proved: for
$A\in U_{\mathrm{AF},j}$ (in the app229 sense) and horizontal $X$,
$\|X\|_{L^{4}}^{2}\le c_{\mathrm{horiz}}\,(j+1)/2^{j}\,
\|D_{A}X\|_{L^{2}}^{2}$.
\item \textbf{The distinction between $M_{A}^{\mathrm{FP}}$ on
$0$-forms and $M_{A}^{\mathrm{YM}}$ on $1$-forms}
(Remark~\ref{rem:two_hessians}) is correct: the Gribov obstruction
of app228 concerns $M_{A}^{\mathrm{FP}}$, not $M_{A}^{\mathrm{YM}}$.
\item \textbf{The gauge-invariant perspective} (replacing
$\Phi_{j}=-\tfrac{1}{2}\log\Delta_{\mathrm{FP},j}$ with
$\Psi_{j}=\tfrac{1}{2}S_{\mathrm{YM},j}$) is the correct
framework, even though the execution has errors.
\end{enumerate}
\end{remark}

\subsection{The corrected bracket bound}
\label{ssec:app230_fix}

The error in Lemma~\ref{lem:bracket_bound} (Theorem~\ref{thm:app229_error1})
is \emph{correctable} by redefining the AF region using the
\emph{gauge-invariant} norm $\|F_{A}\|_{L^{2}}$ instead of
$\|A\|_{L^{2}}$.

\begin{definition}[Corrected AF region]
\label{def:af_region_corr}
The \emph{corrected asymptotic-freedom region} at scale $j$ is
\begin{equation}
\label{eq:af_region_corr}
  U_{\mathrm{AF},j}^{\mathrm{corr}}
  \;:=\;\Bigl\{A\in\mathcal{A}_{j}:
  \|F_{A}\|_{L^{2}(\mathbb{T}^{4})}\le
  \eta_{\mathrm{AF}}^{\mathrm{corr}}\,\frac{2^{2j}}{j+1}\Bigr\},
\end{equation}
where $\eta_{\mathrm{AF}}^{\mathrm{corr}}>0$ is a constant to be
chosen.
\end{definition}

\begin{remark}[Comparison with the app229 AF region]
\label{rem:af_comparison}
The app229 AF region
$U_{\mathrm{AF},j}=\{\|A\|_{L^{2}}\le\eta_{\mathrm{AF}}\,
2^{j}/(j+1)^{1/2}\}$ controls $\|A\|_{L^{2}}$ but only gives
$\|F_{A}\|_{L^{2}}\le 2\eta_{\mathrm{AF}}\,2^{2j}/(j+1)^{1/2}$
(Equation~\eqref{eq:fa_bound_correct}), which is weaker by
$(j+1)^{1/2}$ than needed. The corrected region
$U_{\mathrm{AF},j}^{\mathrm{corr}}$ directly controls
$\|F_{A}\|_{L^{2}}\le\eta_{\mathrm{AF}}^{\mathrm{corr}}\,
2^{2j}/(j+1)$, which is the bound needed for the bracket estimate.

The corrected region is \emph{smaller} than the app229 region:
$U_{\mathrm{AF},j}^{\mathrm{corr}}\subset U_{\mathrm{AF},j}$ for
$\eta_{\mathrm{AF}}^{\mathrm{corr}}$ chosen appropriately (since
$\|F_{A}\|_{L^{2}}\ge\|dA\|_{L^{2}}-\tfrac{1}{2}\|[A,A]\|_{L^{
\mathrm{2}}}\ge 2^{j}\|A\|_{L^{2}}-c\|A\|_{L^{2}}^{2}$, and on
$U_{\mathrm{AF},j}^{\mathrm{corr}}$ this gives
$\|A\|_{L^{2}}\le(\eta_{\mathrm{AF}}^{\mathrm{corr}}\,2^{2j}/(j+1)
+c\|A\|_{L^{2}}^{2})/2^{j}$, hence
$\|A\|_{L^{2}}\le C\,\eta_{\mathrm{AF}}^{\mathrm{corr}}\,
2^{j}/(j+1)$ for small $\eta_{\mathrm{AF}}^{\mathrm{corr}}$).
\end{remark}

\begin{lemma}[Corrected bracket bound]
\label{lem:bracket_bound_corr}
There exist constants $C_{\mathrm{br}},\eta_{\mathrm{AF}}^{
\mathrm{corr}}>0$, independent of $j$ and $n$, such that for all
$j\ge 0$, all $A\in U_{\mathrm{AF},j}^{\mathrm{corr}}$, and all
horizontal $X\in\mathcal{H}_{A}$ at scale $j$,
\begin{equation}
\label{eq:bracket_bound_corr}
  \bigl|\langle F_{A},[X\wedge X]\rangle_{L^{2}}\bigr|
  \;\le\;C_{\mathrm{br}}\,\eta_{\mathrm{AF}}^{\mathrm{corr}}\,
  \|D_{A}X\|_{L^{2}}^{2}.
\end{equation}
\end{lemma}

\begin{proof}
By the scale-$j$ Sobolev inequality
\eqref{eq:scale_j_sobolev_recap},
$\|X\|_{L^{4}}^{2}\le c_{\mathrm{Sob}}\,(j+1)^{1/2}/2^{j}\,
\|X\|_{L^{2}}\,\|dX\|_{L^{2}}$. Using the horizontal Poincar\'e
$\|X\|_{L^{2}}\le(j+1)^{1/2}/2^{j}\,\|D_{A}X\|_{L^{2}}$
(Lemma~\ref{lem:horizontal_sobolev}, applied at scale $j$; this
lemma holds on $U_{\mathrm{AF},j}^{\mathrm{corr}}\subset
U_{\mathrm{AF},j}$ by Remark~\ref{rem:af_comparison}, and the
commutator $\|[A,X]\|_{L^{2}}$ is absorbed as in the proof of
Lemma~\ref{lem:horizontal_sobolev}):
\begin{equation}
\label{eq:x_l4_corrected}
  \|X\|_{L^{4}}^{2}
  \;\le\;c_{\mathrm{Sob}}\,\frac{(j+1)^{1/2}}{2^{j}}\,
  \frac{(j+1)^{1/2}}{2^{j}}\,\|D_{A}X\|_{L^{2}}^{2}
  \;=\;c_{\mathrm{Sob}}\,\frac{j+1}{2^{2j}}\,
  \|D_{A}X\|_{L^{2}}^{2}.
\end{equation}
(Here we used $\|dX\|_{L^{2}}\le\|D_{A}X\|_{L^{2}}+\|[A,X]\|_{
\mathrm{L^{2}}}\le 2\|D_{A}X\|_{L^{2}}$ on $U_{\mathrm{AF},j}^{
\mathrm{corr}}$, absorbing the commutator as in
Lemma~\ref{lem:horizontal_sobolev}.)

By the Cauchy--Schwarz inequality and
$\|[X\wedge X]\|_{L^{2}}\le\|X\|_{L^{4}}^{2}$, substituting
\eqref{eq:x_l4_corrected}:
\begin{equation}
\label{eq:bracket_final_corr}
  \bigl|\langle F_{A},[X\wedge X]\rangle_{L^{2}}\bigr|
  \;\le\;\|F_{A}\|_{L^{2}}\,\|X\|_{L^{4}}^{2}
  \;\le\;\eta_{\mathrm{AF}}^{\mathrm{corr}}\,\frac{2^{2j}}{j+1}
  \;\cdot\;c_{\mathrm{Sob}}\,\frac{j+1}{2^{2j}}\,
  \|D_{A}X\|_{L^{2}}^{2}
  \;=\;c_{\mathrm{Sob}}\,\eta_{\mathrm{AF}}^{\mathrm{corr}}\,
  \|D_{A}X\|_{L^{2}}^{2}.
\end{equation}
This is \eqref{eq:bracket_bound_corr} with
$C_{\mathrm{br}}=c_{\mathrm{Sob}}$.
\end{proof}

\begin{remark}[The key cancellation]
\label{rem:key_cancellation}
The corrected bracket bound \eqref{eq:bracket_final_corr} works
because the scale-$j$ factors cancel exactly:
\begin{equation*}
  \underbrace{\frac{2^{2j}}{j+1}}_{\text{from }\|F_{A}\|_{L^{2}}}
  \;\cdot\;
  \underbrace{\frac{j+1}{2^{2j}}}_{\text{from }\|X\|_{L^{4}}^{2}}
  \;=\;1.
\end{equation*}
This cancellation requires:
\begin{enumerate}
\item The AF region controls $\|F_{A}\|_{L^{2}}\le
\eta_{\mathrm{AF}}^{\mathrm{corr}}\,2^{2j}/(j+1)$ (gauge-invariant).
\item The horizontal Sobolev gives
$\|X\|_{L^{4}}^{2}\le C\,(j+1)/2^{2j}\,\|D_{A}X\|_{L^{2}}^{2}$.
\end{enumerate}
The app229 AF region (controlling $\|A\|_{L^{2}}$) gives
$\|F_{A}\|_{L^{2}}\le 2\eta_{\mathrm{AF}}\,2^{2j}/(j+1)^{1/2}$
(the linear term $dA$ dominates), which is too weak by
$(j+1)^{1/2}$. The corrected AF region (controlling
$\|F_{A}\|_{L^{2}}$ directly) restores the cancellation.
\end{remark}

\begin{theorem}[Corrected Hessian coercivity]
\label{thm:hessian_coercivity_corr}
There exist constants $\kappa_{\mathrm{YM}},\eta_{\mathrm{AF}}^{
\mathrm{corr}}>0$, independent of $j$ and $n$, such that for all
$j\ge 0$, all $A\in U_{\mathrm{AF},j}^{\mathrm{corr}}$, and all
horizontal $X\in\mathcal{H}_{A}$ at scale $j$,
\begin{equation}
\label{eq:hessian_coercivity_corr}
  \mathrm{Hess}\,S_{\mathrm{YM},j}(A)(X,X)
  \;\ge\;\kappa_{\mathrm{YM}}\,\frac{2^{2j}}{j+1}\,
  \|X\|_{L^{2}}^{2}.
\end{equation}
\end{theorem}

\begin{proof}
By the Hessian formula \eqref{eq:ym_hessian_horizontal} and the
corrected bracket bound (Lemma~\ref{lem:bracket_bound_corr}):
\begin{equation*}
  \mathrm{Hess}\,S_{\mathrm{YM},j}(A)(X,X)
  \;=\;\frac{1}{g_{\mathrm{YM}}^{2}}\,
  \bigl(\|D_{A}X\|_{L^{2}}^{2}
  +\langle F_{A},[X\wedge X]\rangle_{L^{2}}\bigr)
  \;\ge\;\frac{1}{g_{\mathrm{YM}}^{2}}\,
  \bigl(1-C_{\mathrm{br}}\,\eta_{\mathrm{AF}}^{\mathrm{corr}}\bigr)\,
  \|D_{A}X\|_{L^{2}}^{2}.
\end{equation*}
Choosing $\eta_{\mathrm{AF}}^{\mathrm{corr}}\le 1/(2C_{\mathrm{br}})$:
\begin{equation}
\label{eq:hessian_intermediate_corr}
  \mathrm{Hess}\,S_{\mathrm{YM},j}(A)(X,X)
  \;\ge\;\frac{1}{2g_{\mathrm{YM}}^{2}}\,\|D_{A}X\|_{L^{2}}^{2}.
\end{equation}
By the horizontal Sobolev inequality
(Lemma~\ref{lem:horizontal_sobolev}, which holds on
$U_{\mathrm{AF},j}^{\mathrm{corr}}\subset U_{\mathrm{AF},j}$):
$\|D_{A}X\|_{L^{2}}^{2}\ge c_{\mathrm{horiz}}\,2^{2j}/(j+1)\,
\|X\|_{L^{2}}^{2}$. Substituting into
\eqref{eq:hessian_intermediate_corr}:
\begin{equation*}
  \mathrm{Hess}\,S_{\mathrm{YM},j}(A)(X,X)
  \;\ge\;\frac{c_{\mathrm{horiz}}}{2g_{\mathrm{YM}}^{2}}\,
  \frac{2^{2j}}{j+1}\,\|X\|_{L^{2}}^{2},
\end{equation*}
which is \eqref{eq:hessian_coercivity_corr} with
$\kappa_{\mathrm{YM}}=c_{\mathrm{horiz}}/(2g_{\mathrm{YM}}^{2})$.
\end{proof}

\begin{theorem}[Corrected BE curvature]
\label{thm:be_curvature_corr}
There exists $\kappa_{0}^{\mathrm{corr}}>0$, independent of $j$
and $n$, such that for all $j\ge 0$ and all
$A\in U_{\mathrm{AF},j}^{\mathrm{corr}}$,
\begin{equation}
\label{eq:be_curvature_corr}
  \bigl(\mathrm{Ric}_{g_{j}}+\mathrm{Hess}_{g_{j}}\,\Psi_{j}
  \bigr)(A)
  \;\ge\;\kappa_{0}^{\mathrm{corr}}\,\frac{2^{2j}}{j+1}\,g_{j}(A),
\end{equation}
where $\Psi_{j}=\tfrac{1}{2}S_{\mathrm{YM},j}-\log Z_{j}$.
\end{theorem}

\begin{proof}
Identical to the proof of
Theorem~\ref{thm:be_curvature_gauge_invariant} (app229), but with
the corrected Hessian coercivity
(Theorem~\ref{thm:hessian_coercivity_corr}) replacing the
erroneous app229 version. The BE curvature is
$\mathrm{Ric}_{g_{j}}+\mathrm{Hess}\,\Psi_{j}$, and
$\mathrm{Hess}\,\Psi_{j}=\tfrac{1}{2}\mathrm{Hess}\,S_{\mathrm{YM},j}$
(since $\log Z_{j}$ is constant in $A$). The metric
$g_{j}$ on $\mathcal{A}_{j}$ is flat, so $\mathrm{Ric}_{g_{j}}=0$.
By Theorem~\ref{thm:hessian_coercivity_corr}:
\begin{equation*}
  \mathrm{Hess}\,\Psi_{j}(A)(X,X)
  \;=\;\tfrac{1}{2}\,\mathrm{Hess}\,S_{\mathrm{YM},j}(A)(X,X)
  \;\ge\;\frac{\kappa_{\mathrm{YM}}}{2}\,\frac{2^{2j}}{j+1}\,
  \|X\|_{L^{2}}^{2}
  \;=\;\frac{\kappa_{\mathrm{YM}}}{2}\,\frac{2^{2j}}{j+1}\,
  g_{j}(X,X),
\end{equation*}
which is \eqref{eq:be_curvature_corr} with
$\kappa_{0}^{\mathrm{corr}}=\kappa_{\mathrm{YM}}/2$.
\end{proof}

\begin{theorem}[Corrected per-scale Cheeger constant]
\label{thm:sharp_cheeger_corr}
There exists $c_{0}^{\mathrm{corr}}>0$, independent of $j$ and
$n$, such that for all $j\ge 0$,
\begin{equation}
\label{eq:sharp_cheeger_corr}
  h_{j}^{\nu}\bigl(U_{\mathrm{AF},j}^{\mathrm{corr}}\bigr)
  \;\ge\;c_{0}^{\mathrm{corr}}\,\frac{2^{j}}{(j+1)^{1/2}}.
\end{equation}
\end{theorem}

\begin{proof}
By the Bakry--\'Emery Cheeger bound
(Bakry, Gentil \& Ledoux 2014, Theorem~4.7.2) and the corrected
BE curvature (Theorem~\ref{thm:be_curvature_corr}):
\begin{equation*}
  h_{j}^{\nu}\bigl(U_{\mathrm{AF},j}^{\mathrm{corr}}\bigr)
  \;\ge\;\frac{\sqrt{\kappa_{0}^{\mathrm{corr}}\,2^{2j}/(j+1)}}{2}
  \;=\;\frac{\sqrt{\kappa_{0}^{\mathrm{corr}}}}{2}\,
  \frac{2^{j}}{(j+1)^{1/2}},
\end{equation*}
which is \eqref{eq:sharp_cheeger_corr} with
$c_{0}^{\mathrm{corr}}=\sqrt{\kappa_{0}^{\mathrm{corr}}}/2$.
\end{proof}

\begin{remark}[The corrected bound is still at $\epsilon=1/2$]
\label{rem:corrected_still_half}
The corrected per-scale Cheeger bound
$h_{j}^{\nu}\ge c_{0}^{\mathrm{corr}}\,2^{j}/(j+1)^{1/2}$
(Theorem~\ref{thm:sharp_cheeger_corr}) is the \emph{same} as the
app229 bound (Theorem~\ref{thm:sharp_cheeger_gauge_invariant}),
but now with a \emph{correct} proof. In the notation of app226
(Theorem~\ref{thm:sharp_per_scale_cheeger}), this corresponds to
$\epsilon=1/2$, the boundary case.

However, as shown in Theorem~\ref{thm:app229_error2}, the
multiscale aggregation of the per-scale bound
$h_{j}^{\nu}\sim 2^{j}/(j+1)^{1/2}$ gives
$h_{\mathrm{OS}}^{(n)}\sim c/L_{\mathrm{phys}}$ (in physical
units), which corresponds to $h_{\mathrm{RP}}^{(n)}\sim c$ (in
lattice units), \emph{not} $h_{\mathrm{RP}}^{(n)}\sim c\sqrt{a_{
\mathrm{n}}}$. The mass gap requires the latter.
\end{remark}

\subsection{The remaining obstruction: gauge-invariant concentration}
\label{ssec:app230_concentration}

The corrected bracket bound (Theorem~\ref{thm:hessian_coercivity_corr})
and BE curvature (Theorem~\ref{thm:be_curvature_corr}) hold on
$U_{\mathrm{AF},j}^{\mathrm{corr}}$. To extend the Cheeger bound
from $U_{\mathrm{AF},j}^{\mathrm{corr}}$ to all of $\mathcal{Q}_{j}$
(as in Step 3 of the proof of
Theorem~\ref{thm:sharp_cheeger_gauge_invariant}), one needs the
concentration estimate:
\begin{equation}
\label{eq:concentration_needed}
  \nu_{j}\bigl(\mathcal{Q}_{j}\setminus
  U_{\mathrm{AF},j}^{\mathrm{corr}}\bigr)
  \;\le\;c_{\mathrm{AF}}\,\exp\!\bigl(-\alpha_{\mathrm{AF}}\,
  2^{2j}/(j+1)\bigr).
\end{equation}

\begin{hypothesis}[Gauge-invariant concentration]
\label{hyp:gauge_invariant_concentration}
There exist constants $c_{\mathrm{AF}},\alpha_{\mathrm{AF}}>0$,
independent of $j$ and $n$, such that
\eqref{eq:concentration_needed} holds for all $j\ge 0$, where
$\nu_{j}$ is the gauge-invariant quotient measure
(Definition~\ref{def:gauge_invariant_measure}) and
$U_{\mathrm{AF},j}^{\mathrm{corr}}$ is the corrected AF region
(Definition~\ref{def:af_region_corr}).
\end{hypothesis}

\begin{remark}[Why the concentration is non-trivial]
\label{rem:concentration_hard}
The concentration estimate \eqref{eq:concentration_needed} for the
corrected AF region $U_{\mathrm{AF},j}^{\mathrm{corr}}=\{\|F_{A}\|_{
\mathrm{L^{2}}}\le\eta_{\mathrm{AF}}^{\mathrm{corr}}\,2^{2j}/(j+1)\}$
is \emph{harder} than the app229 concentration for
$U_{\mathrm{AF},j}=\{\|A\|_{L^{2}}\le\eta_{\mathrm{AF}}\,2^{j}/
(j+1)^{1/2}\}$, because:
\begin{enumerate}
\item The corrected region is \emph{smaller} (it controls the
gauge-invariant $\|F_{A}\|_{L^{2}}$, not just $\|A\|_{L^{2}}$).
\item The complement
$\{\|F_{A}\|_{L^{2}}>\eta_{\mathrm{AF}}^{\mathrm{corr}}\,2^{2j}/
(j+1)\}$ includes connections with small $\|A\|_{L^{2}}$ but
large $\|F_{A}\|_{L^{2}}$ (e.g.\ near-instantons with small
radius), whose $\nu_{j}$-measure is not obviously exponentially
small.
\item The app229 concentration proof (Lemma~\ref{lem:af_concentration})
uses the lower bound $S_{\mathrm{YM},j}(A)\ge\|dA\|_{L^{2}}^{2}/
(2g_{\mathrm{YM}}^{2})\ge 2^{2j}\|A\|_{L^{2}}^{2}/(2g_{\mathrm{YM}}^{
\mathrm{2}})$, which controls $\|A\|_{L^{2}}$ but \emph{not}
$\|F_{A}\|_{L^{2}}$ directly (due to possible cancellation between
$dA$ and $[A,A]$).
\end{enumerate}

A proof of Hypothesis~\ref{hyp:gauge_invariant_concentration}
would require either:
\begin{itemize}
\item A \emph{gauge-invariant} lower bound on the YM action of the
form $S_{\mathrm{YM},j}(A)\ge c\,\|F_{A}\|_{L^{2}}^{2}/g_{\mathrm{YM}}^{
\mathrm{2}}$ (which holds, since $S_{\mathrm{YM}}=\|F_{A}\|_{L^{
\mathrm{2}}}^{2}/(2g_{\mathrm{YM}}^{2})$ by definition), combined
with a \emph{partition-function lower bound} that controls the
denominator $Z_{j}$, or
\item A \emph{Ward-identity argument} that controls the
$\nu_{j}$-measure of connections with large $\|F_{A}\|_{L^{2}}$
relative to $\|A\|_{L^{2}}$ (i.e.\ connections near the Gribov
horizon, where the gauge slice degenerates).
\end{itemize}
The first approach is straightforward: on
$U_{\mathrm{AF},j}^{\mathrm{corr},c}=\{\|F_{A}\|_{L^{2}}>\eta_{
\mathrm{AF}}^{\mathrm{corr}}\,2^{2j}/(j+1)\}$,
$S_{\mathrm{YM},j}(A)=\|F_{A}\|_{L^{2}}^{2}/(2g_{\mathrm{YM}}^{2})>
(\eta_{\mathrm{AF}}^{\mathrm{corr}})^{2}\,2^{4j}/(2g_{\mathrm{YM}}^{
\mathrm{2}}(j+1)^{2})$. The Chebyshev bound gives
\begin{equation}
\label{eq:chebyshev_gauge_invariant}
  \nu_{j}(U_{\mathrm{AF},j}^{\mathrm{corr},c})
  \;\le\;\exp\!\bigl(-(\eta_{\mathrm{AF}}^{\mathrm{corr}})^{2}\,
  2^{4j}/(2g_{\mathrm{YM}}^{2}(j+1)^{2})\bigr)\,
  \frac{\mathrm{vol}_{g_{j}}(\mathcal{Q}_{j})}{Z_{j}}.
\end{equation}
The exponential rate is
$\sim 2^{4j}/(g_{\mathrm{YM}}^{2}(j+1)^{2})$, which by asymptotic
freedom ($1/g_{\mathrm{YM}}^{2}\sim\log(2^{j}/\Lambda_{\mathrm{QCD}})$)
is $\sim 2^{4j}\log(2^{j})/(j+1)^{2}$, much \emph{larger} than the
required $2^{2j}/(j+1)$. So
\eqref{eq:concentration_needed} holds \emph{with room to spare},
\emph{provided} the volume-to-partition-function ratio
$\mathrm{vol}_{g_{j}}(\mathcal{Q}_{j})/Z_{j}$ is bounded.
\end{remark}

\begin{lemma}[Volume-to-partition-function ratio]
\label{lem:vol_ratio}
The ratio $\mathrm{vol}_{g_{j}}(\mathcal{Q}_{j})/Z_{j}$ is bounded
by a constant independent of $j$ and $n$.
\end{lemma}

\begin{proof}[Sketch]
The partition function $Z_{j}=\int_{\mathcal{Q}_{j}}
e^{-S_{\mathrm{YM},j}}\,d\mathrm{vol}_{g_{j}}$ is bounded below
by the contribution from a neighborhood of the minimum of
$S_{\mathrm{YM},j}$ (the flat connection $A=0$, where
$S_{\mathrm{YM},j}(0)=0$):
\begin{equation*}
  Z_{j}
  \;\ge\;\int_{\|F_{A}\|_{L^{2}}\le\eta_{\mathrm{AF}}^{\mathrm{corr}}\,
  2^{2j}/(j+1)}
  e^{-S_{\mathrm{YM},j}}\,d\mathrm{vol}_{g_{j}}
  \;\ge\;e^{-\alpha_{\mathrm{AF}}\,2^{2j}/(j+1)}\,
  \mathrm{vol}_{g_{j}}(U_{\mathrm{AF},j}^{\mathrm{corr}}),
\end{equation*}
using $S_{\mathrm{YM},j}(A)\le(\eta_{\mathrm{AF}}^{\mathrm{corr}})^{
\mathrm{2}}\,2^{2j}/(2g_{\mathrm{YM}}^{2}(j+1))$ on
$U_{\mathrm{AF},j}^{\mathrm{corr}}$. If
$\mathrm{vol}_{g_{j}}(U_{\mathrm{AF},j}^{\mathrm{corr}})\ge
c\,\mathrm{vol}_{g_{j}}(\mathcal{Q}_{j})$ for a constant $c>0$
(which requires that $U_{\mathrm{AF},j}^{\mathrm{corr}}$ has a
uniform fraction of the total volume --- a non-trivial geometric
estimate), then
$\mathrm{vol}_{g_{j}}(\mathcal{Q}_{j})/Z_{j}\le
c^{-1}\,e^{\alpha_{\mathrm{AF}}\,2^{2j}/(j+1)}$, and
\eqref{eq:chebyshev_gauge_invariant} gives
\eqref{eq:concentration_needed} (the exponential rate
$2^{4j}/(g_{\mathrm{YM}}^{2}(j+1)^{2})$ dominates
$\alpha_{\mathrm{AF}}\,2^{2j}/(j+1)$ for large $j$).

The volume estimate
$\mathrm{vol}_{g_{j}}(U_{\mathrm{AF},j}^{\mathrm{corr}})\ge
c\,\mathrm{vol}_{g_{j}}(\mathcal{Q}_{j})$ is the non-trivial
input; it is equivalent to
$\nu_{j}(U_{\mathrm{AF},j}^{\mathrm{corr}})\ge c'>0$ for a
constant $c'>0$, which is a weaker form of
Hypothesis~\ref{hyp:gauge_invariant_concentration} (without the
exponential rate). This circular dependency means the
concentration must be proved by a different method.
\end{proof}

\begin{remark}[The circular dependency]
\label{rem:circular}
The Chebyshev argument (Remark~\ref{rem:concentration_hard}) reduces
the concentration \eqref{eq:concentration_needed} to the
volume-to-partition-function ratio
$\mathrm{vol}_{g_{j}}(\mathcal{Q}_{j})/Z_{j}$, which in turn
requires a lower bound on $Z_{j}$ in terms of
$\mathrm{vol}_{g_{j}}(U_{\mathrm{AF},j}^{\mathrm{corr}})$, which
requires a lower bound on
$\mathrm{vol}_{g_{j}}(U_{\mathrm{AF},j}^{\mathrm{corr}})$, which
is a weaker form of the concentration. This circular dependency
is the obstruction.

A way out is to prove the concentration by a \emph{different}
method, e.g.:
\begin{enumerate}
\item \textbf{BRST/Slavnov--Taylor identity}: the gauge-invariant
measure $\nu_{j}$ satisfies a Ward identity that controls the
distribution of $\|F_{A}\|_{L^{2}}$ directly.
\item \textbf{Lattice a-priori estimate}: on the lattice, the
Wilson action $S_{\mathrm{Wil}}\ge c\sum_{p}|1-U_{p}|^{2}\ge
c'\|F_{A}\|_{L^{2}}^{2}$ (for weak fields), and the lattice
partition function can be bounded by a Gaussian integral.
\item \textbf{Cluster expansion}: for large $\beta$ (weak coupling),
the measure concentrates on weak fields by a cluster expansion
(Balaban 1984, Federbush 1984), giving the required
concentration.
\end{enumerate}
Each of these is a substantial technical project.
\end{remark}

\subsection{Honest final status}
\label{ssec:app230_status}

\begin{theorem}[Status after app229--app230]
\label{thm:status_app230}
The Clay Millennium Prize mass-gap conjecture for 4D pure
Yang--Mills remains \textbf{open}. The proof strategy of app229
(gauge-invariant quotient measure, BE curvature, Cheeger
constant) is \emph{partially} correct:
\begin{enumerate}
\item \textbf{Corrected} (Theorem~\ref{thm:hessian_coercivity_corr}):
the YM Hessian on horizontal $1$-forms is uniformly coercive on
the corrected AF region
$U_{\mathrm{AF},j}^{\mathrm{corr}}=\{\|F_{A}\|_{L^{2}}\le
\eta_{\mathrm{AF}}^{\mathrm{corr}}\,2^{2j}/(j+1)\}$.
\item \textbf{Corrected} (Theorem~\ref{thm:be_curvature_corr}):
the BE curvature of the gauge-invariant measure $\nu_{j}$ is
$\ge\kappa_{0}^{\mathrm{corr}}\,2^{2j}/(j+1)$ on
$U_{\mathrm{AF},j}^{\mathrm{corr}}$.
\item \textbf{Corrected} (Theorem~\ref{thm:sharp_cheeger_corr}):
the per-scale Cheeger constant on $U_{\mathrm{AF},j}^{\mathrm{corr}}$
is $\ge c_{0}^{\mathrm{corr}}\,2^{j}/(j+1)^{1/2}$.
\item \textbf{Open} (Hypothesis~\ref{hyp:gauge_invariant_concentration}):
the concentration of $\nu_{j}$ on $U_{\mathrm{AF},j}^{\mathrm{corr}}$
is not proved.
\item \textbf{Open} (Theorem~\ref{thm:app229_error2}): the
scaling/normalization relating the per-scale Cheeger constant
$h_{j}^{\nu}$ (physical units) to the OS transfer-matrix spectral
gap $1-\lambda_{1}^{(n)}$ (lattice units) requires an additional
factor of $a_{n}^{1/2}$ that is not provided by the current
argument.
\end{enumerate}
The mass gap is reduced to:
\begin{itemize}
\item Hypothesis~\ref{hyp:gauge_invariant_concentration}
(gauge-invariant concentration), \emph{and}
\item a correct normalization of the multiscale aggregation that
accounts for the $a_{n}$-dependence of the OS transfer-matrix
Cheeger constant.
\end{itemize}
\end{theorem}

\begin{remark}[Comparison with app226--app228]
\label{rem:comparison_app226}
The status after app230 is \emph{essentially the same} as after
app228: the mass gap is reduced to a single sharp input. The
difference is that the input is now \emph{gauge-invariant}
(Hypothesis~\ref{hyp:gauge_invariant_concentration}, controlling
$\|F_{A}\|_{L^{2}}$) rather than \emph{gauge-fixed}
(app228, controlling the Faddeev--Popov operator). The
gauge-invariant formulation avoids the Gribov obstruction of
app228, but the concentration estimate is correspondingly harder.

The shortest credible route to the Clay mass gap remains:
\begin{enumerate}
\item Prove Hypothesis~\ref{hyp:gauge_invariant_concentration}
(gauge-invariant concentration).
\item Correct the scaling/normalization of the multiscale
aggregation (Theorem~\ref{thm:app229_error2}), or alternatively
prove the per-scale Cheeger bound at
$h_{j}^{\nu}\ge c\,2^{2j}/(j+1)$ (which would give the correct
$a_{n}^{1/2}$ scaling automatically).
\end{enumerate}
\end{remark}

\begin{remark}[New results proved in this appendix]
\label{rem:app230_new}
This appendix proves the following new results:
\begin{enumerate}
\item \textbf{Theorem~\ref{thm:app229_error1}}: identification of
the bracket bound error in app229.
\item \textbf{Theorem~\ref{thm:app229_error2}}: identification of
the scaling/normalization error in app229.
\item \textbf{Lemma~\ref{lem:bracket_bound_corr}}: corrected
bracket bound using the gauge-invariant AF region.
\item \textbf{Theorem~\ref{thm:hessian_coercivity_corr}}:
corrected Hessian coercivity on the gauge-invariant AF region.
\item \textbf{Theorem~\ref{thm:be_curvature_corr}}: corrected
BE curvature bound.
\item \textbf{Theorem~\ref{thm:sharp_cheeger_corr}}: corrected
per-scale Cheeger constant.
\end{enumerate}
\end{remark}

\endinput
